\documentclass[../../../include/open-logic-section]{subfiles}

\begin{document}
\olfileid[hi]{sth}{card-arithmetic}{simp}

\olsection{योग और गुणन का सरलीकरण}

पता चलता है कि परिमितेतर कार्डिनल योग और गुणन \emph{अत्यंत} सरल हैं। यह इस
तथ्य से निष्पन्न होता है कि कार्डिनल (कुछ) क्रमसूचक हैं, अतः सुव्यवस्थित हैं
और इसलिए एक विशेष ढंग से उनका परिचालन किया जा सकता है। किंतु इसे दिखाना
इतना सरल \emph{नहीं} है। आरंभ में हमें एक युक्तिपूर्ण परिभाषा चाहिए:

\begin{defn}
हम क्रमसूचकों के युग्मों पर एक \emph{विहित क्रम} $\canonord$ इस शर्त से
परिभाषित करते हैं कि $\tuple{\alpha_1,\alpha_2}\canonord
\tuple{\beta_1,\beta_2}$ तभी जब इनमें से कोई एक हो:
\begin{enumerate}
	\item $\max(\alpha_1, \alpha_2)<\max(\beta_1, \beta_2)$; अथवा
	\item $\max(\alpha_1, \alpha_2) = \max(\beta_1, \beta_2)$ और
	$\alpha_1 < \beta_1$; अथवा
	\item $\max(\alpha_1, \alpha_2) = \max(\beta_1, \beta_2)$ और
	$\alpha_1 = \beta_1$ और $\alpha_2 < \beta_2$
\end{enumerate}
\end{defn}

\begin{lem}
किसी भी क्रमसूचक $\alpha$ के लिए
$\tuple{\alpha\times\alpha,\canonord}$ एक सुव्यवस्था है।
\end{lem}

\begin{proof}
स्पष्टतः $\canonord$, $\alpha\times\alpha$ पर संबद्ध है। मान लें कि न
$\tuple{\alpha_1,\alpha_2}$, $\tuple{\beta_1,\beta_2}$ से
$\canonord$-लघुतर है, न दूसरा पहले से। तब
$\max(\alpha_1,\alpha_2)=\max(\beta_1,\beta_2)$ और
$\alpha_1=\beta_1$ और $\alpha_2=\beta_2$, इसलिए
$\tuple{\alpha_1,\alpha_2}=\tuple{\beta_1,\beta_2}$।

सुव्यवस्था दिखाने के लिए $X\subseteq\alpha\times\alpha$ को अरिक्त मानें।
चूँकि $\alpha$ क्रमसूचक है, कोई $\delta$ समुच्चय
$\Setabs{\max(\gamma_1,\gamma_2)}{\tuple{\gamma_1,\gamma_2}\in X}$ का
लघुतम सदस्य है। अब
$\Setabs{\tuple{\gamma_1,\gamma_2}\in X}{\max(\gamma_1,\gamma_2)=\delta}$
से उन युग्मों को छोड़ दें जिनका पहला निर्देशांक लघुतम नहीं है; शेष युग्मों
में लघुतम दूसरे निर्देशांक वाला युग्म $X$ का $\canonord$-लघुतम अवयव है।
\end{proof}
\noindent
अब एक छोटा-सा सरल प्रेक्षण:

\begin{prop}\ollabel{simplecardproduct}
यदि $\cardeq{\alpha}{\beta}$, तो
$\cardeq{\alpha\times\alpha}{\beta\times\beta}$।
\end{prop}

\begin{proof}
बस $f\colon\alpha\to\beta$ से
$\tuple{\gamma_1,\gamma_2}\mapsto\tuple{f(\gamma_1),f(\gamma_2)}$
प्रेरित होने दें।
\end{proof}
\noindent
अब हम इन सबका उपयोग एक महत्त्वपूर्ण लेम्मा सिद्ध करने में करेंगे:
\begin{lem}\ollabel{alphatimesalpha}
किसी भी अनंत क्रमसूचक $\alpha$ के लिए
$\cardeq{\alpha}{\alpha\times\alpha}$।
\end{lem}

\begin{proof}
विरोधाभास द्वारा $\alpha$ को वह लघुतम अनंत क्रमसूचक मानें जिसके लिए यह
असत्य है। \olref[sfr][siz][zigzag]{natsquaredenumerable} दिखाता है कि
$\cardeq{\omega}{\omega\times\omega}$, अतः $\omega\in\alpha$। इसके
अतिरिक्त $\alpha$ एक कार्डिनल है: विरोधाभास के लिए अन्यथा मान लें; तब
$\card{\alpha}\in\alpha$, इसलिए परिकल्पना से
$\cardeq{\card{\alpha}}{\card{\alpha}\times\card{\alpha}}$; और परिभाषा से
$\cardeq{\card{\alpha}}{\alpha}$; अतः \olref{simplecardproduct} से
$\cardeq{\alpha}{\alpha\times\alpha}$।

अब प्रत्येक $\tuple{\gamma_1,\gamma_2}\in\alpha\times\alpha$ के लिए इस
खंड पर विचार करें:
\begin{align*}
	\text{Seg}(\gamma_1, \gamma_2) &= \Setabs{\tuple{\delta_1, \delta_2} \in \alpha \times \alpha}{\tuple{\delta_1, \delta_2} \canonord \tuple{\gamma_1, \gamma_2}}
\end{align*}
$\gamma=\max(\gamma_1,\gamma_2)$ रखते हुए ध्यान दें कि
$\tuple{\gamma_1,\gamma_2}\canonord\tuple{\gamma+1,\gamma+1}$। अतः जब
$\gamma$ अनंत हो, ध्यान दें:
\begin{align*}
	\text{Seg}(\gamma_1, \gamma_2) & 
	\precsim ((\gamma \ordplus 1)\ordtimes (\gamma \ordplus 1))\\
	&\approx (\gamma \ordtimes \gamma)
	\text{, \olref[ord-arithmetic][using-addition]{ordinfinitycharacter} और
	\olref{simplecardproduct} से}\\
	&\approx \gamma \text{, आगमन परिकल्पना से}\\
	& \prec \alpha\text{, क्योंकि $\alpha$ एक कार्डिनल है}
\end{align*}
अतः $\ordtype{\alpha\times\alpha,\canonord}\leq\alpha$, और फलतः
$\cardle{\alpha\times\alpha}{\alpha}$। चूँकि निस्संदेह
$\cardle{\alpha}{\alpha\times\alpha}$, परिणाम श्रेडर--बर्नस्टीन से निष्पन्न
होता है।
\end{proof}

अंततः हम अपने सरलीकरण परिणाम पर पहुँचते हैं:

\begin{thm}\ollabel{cardplustimesmax}
यदि $\cardfont{a},\cardfont{b}$ अनंत कार्डिनल हैं, तो:
\[
	\cardfont{a}
\cardtimes \cardfont{b} = \cardfont{a} \cardplus \cardfont{b} =
\text{max}(\cardfont{a}, \cardfont{b}).
\]
\end{thm}

\begin{proof}
व्यापकता खोए बिना मान लें $\cardfont{a}=\max(\cardfont{a},
\cardfont{b})$। तब \olref{alphatimesalpha} का उपयोग करके,
$\cardfont{a}\cardtimes\cardfont{a} = \cardfont{a} \leq \cardfont{a}
\cardplus \cardfont{b} \leq \cardfont{a} \cardplus \cardfont{a} \leq
\cardfont{a}\cardtimes\cardfont{a}$। \end{proof}\noindent इसी प्रकार,
यदि $\cardfont{a}$ अनंत है, तो $\leq\cardfont{a}$ आकार वाले समुच्चयों के
$\cardfont{a}$-आकार वाले सम्मिलन का आकार $\leq\cardfont{a}$ है:

\begin{prop}\ollabel{kappaunionkappasize}
$\cardfont{a}$ को एक अनंत कार्डिनल मानें। प्रत्येक क्रमसूचक
$\beta\in\cardfont{a}$ के लिए $X_\beta$ ऐसा समुच्चय हो कि
$\card{X_\beta}\leq\cardfont{a}$। तब
$\card{\bigcup_{\beta\in\cardfont{a}}X_\beta}\leq\cardfont{a}$।
\end{prop}

\begin{proof}
प्रत्येक $\beta\in\cardfont{a}$ के लिए कोई !!a{injection}
$f_\beta\colon X_\beta\to\cardfont{a}$ नियत करें।\footnote{इन्हें
``नियत'' कैसे किया गया है? \olref[sth][choice][countablechoice]{sec}
देखें।} !!a{injection}
$g\colon\bigcup_{\beta\in\cardfont{a}}X_\beta\to
\cardfont{a}\times\cardfont{a}$ को
$g(v)=\tuple{\beta,f_\beta(v)}$ से परिभाषित करें, जहाँ $v\in X_\beta$ और
किसी $\gamma\in\beta$ के लिए $v\notin X_\gamma$। अब
\olref{cardplustimesmax} से
$\bigcup_{\beta\in\cardfont{a}}X_\beta\preceq
\cardfont{a}\times\cardfont{a}\approx\cardfont{a}$।
\end{proof}

\end{document}
